\chapter{SUMMARY, CONCLUSION, AND RECOMMENDATIONS}

\section{Introduction}

This chapter summarises the results of the research project, the extent to which the identified research objectives were met and the main contributions of the study. It offers a comprehensive overview of the results for the key functional blocks, generalises conclusions on the integration of blockchain technology and predictive analytics for agriculture and proposes concrete suggestions for further technical developments and practical applications.

\section{Summary of Findings}

The task for this project was to develop and deploy the agricultural data management system based on blockchain technology, which is safe and analytics-oriented. The work addressed the core requirements across design, development and evaluation as detailed in the previous chapters.

For the web-based data management platform, the web interface was designed and implemented using a modular approach with React and TypeScript as the front-end framework and Django REST Framework as the back-end framework. The platform validates data on the environmental parameters including temperature in Celsius, Soil moisture percentage, Soil pH, etc through the validated input forms. The data is registered in a PostgreSQL database and displayed in a role-specific dashboard, depending on the needs of the farmers, logistics handlers and system administrators. Both credential-based and Sui wallet authentication options are available for the interface, catering to users who are more or less familiar with blockchain.

In terms of the predictive analytics engine, an engine was created and placed in the back-end. It uses a weighted moving average model from historical telemetry data to estimate yields given the mean value of the historic data for temperature, soil moisture, and soil pH. The prediction formula takes into consideration moisture availability as a positive factor and penalizes watering at temperatures that are not optimum for growth (22 degrees Celsius). A confidence score is calculated depending on the number of collected data points and contextual recommendations are made in order to support the farmer's decision making. The result is stored in Redis to guarantee fast results.

Regarding data integrity and blockchain anchoring, data integrity is enforced through a multi-layered approach. Each telemetry record is secured with a SHA-256 hash that is created from a canonical JSON representation of the core fields of the record. The batch records are created as programmable objects on the Sui blockchain using a Move smart contract that contains the following information: the type of crop, the weight, the farmer of origin, who is currently holding the crops, and the telemetry integrity hash, which is on-chain. The smart contract verifies that transfer can only be executed by the current custodian, thereby enforcing custody transfer between stakeholders. The blockchain anchoring gives rise to an auditable and tamper-proof trace of each blockchain asset correlated with the source of on-chain data.

Finally, on system evaluation and performance, the system was tested in various aspects. All ten functional requirements have been met through the functional testing. Security assessment confirmed the functionality of Argon2id password hashing, JWT token lifecycle control, rate limiting and data integrity hash checking. The performance of the Sui Testnet during blockchain performance testing was found to be 100 percent successful with an average minting latency of 1.2 seconds and transfer latency of 0.9 seconds, along with minimal gas fees. Prediction accuracy was evaluated by consistency and sensitivity analysis, which showed that the model was able to give correct predictions when the input conditions were varied and that the predictions were deterministic.

\section{Conclusion}

The TerraNode project is a tangible example of how to create a comprehensive, agri-led data management ecosystem, seamlessly combining secure data storage, blockchain immutability, supply chain transparency, and predictive analytics in a unified system. Reviews from the literature detailed in Chapter Two showed that these capabilities have been implemented in different ways in different systems, but seldom in combination. By requiring blockchain verification before analytical processing, TerraNode ensures yield predictions are based on data that is mathematically verifiable, bridging the gap between this new data and current data sources.

The decision to use the Sui blockchain for this application seemed appropriate. Sui's delegated proof-of-stake consensus and object-centric data model led to low-cost, sub-second transaction finality for batch minting and custody transfer operations. A linear type system was added to the Move programming language, which offered extra security guarantees at the smart contract level, making it hard to cause common vulnerabilities like double spending of batch objects.

The three-tier architecture, with the React application running on the frontend, the Django application running on the back end, and the Sui blockchain running on a separate layer, helped to ensure a clean separation of concerns, enabling each part to be developed and tested separately. The role based access control system lets each stakeholder type only deal with the functionality relevant to their role, which decreases the intricacy and security dangers.

The current analytics engine is based on a basic weighted moving average model but is intended to support more complex prediction algorithms in the future. With the modular service layer, larger datasets can be used to train machine learning models, adding them to the service without touching the API or the front end.

To conclude, TerraNode is a working prototype that fills an actual need in the agricultural technology sector. It demonstrates that a system that is technically sound and accessible to the smallholder farming communities it's designed for can be created through the integration of modern web technologies with a next-generation blockchain platform and data-driven analytics.

\section{Recommendations for Future Work}

The existing system uses hand entered telemetry data. Future iterations need to be integrated directly with the physical IoT sensors deployed in agriculture fields and the data collected should be able to be done automatically and continuously. This would remove the possibility of data entry errors by people and add more and more frequent data to be analyzed by telemetry.

The weighted moving average model is a useful starting point for yield forecasting, but advanced machine learning models such as Random Forests, Long Short-Term Memory (LSTM) networks, or the CNN-RNN model of \textcite{khaki2019cnn} could greatly enhance prediction accuracy. Such models could also be better generalized if trained on larger, geographically diverse data sets.

At present, the implementation is on the Sui Testnet, which has no real economic value. The deployment of the smart contract on the Sui Mainnet would allow for a true tracking of assets and help integrate the agriculture lending and insurance experience into the decentralized finance (DeFi) ecosystem.

In terms of usability, a mobile phone app for both Android and iOS would allow farmers to access the internet via their mobile phones to be more accessible. The mobile app might also be able to add location and visual crop condition data to the telemetry data by utilizing the device's GPS, camera and other sensors.

The present user interface is in English, and is therefore not accessible in areas where English is not the native language. Support for local languages would greatly increase the potential user base.

With the increasing amount of data on the blockchain, the use of Storage Light Nodes (SLN) proposed by \textcite{sun2025storage} is one of the optimized storage methods in the development of TerraNode, which can help reduce the resource consumption of the nodes involved in the TerraNode network.

Finally, further research is needed to investigate the standardization of agricultural data formats and APIs, facilitating TerraNode's integration with other agricultural information systems, weather services, and market platforms.

\section{Chapter Summary}

This chapter synthesized the findings of the TerraNode project across its four core research objectives, drew conclusions regarding the technical architecture and real-world applicability of blockchain and predictive analytics in agriculture, and proposed seven targeted recommendations for future development. The project demonstrated that combining an object-centric blockchain with deterministic data hashing and responsive web interfaces provides a practical, tamper-resistant data management foundation for smallholder farming communities.